\documentclass[a4paper]{article}
\usepackage{amsfonts}
\usepackage{amsmath}
\usepackage{amssymb}
\usepackage{graphicx}   

\newcommand{\degree}{^{\circ}}
\newcommand{\cis}{\operatorname{cis}}
\newcommand{\C}{\mathbb{C}}

\begin{document}
  
\noindent \large Auteur: Rune Suy \\
\noindent \large Studierichting: Informatica\\
\noindent \large Oefeningenreeks 1C nr. 27.5\\

\medskip

\normalsize

Bereken in $\C$: $ 2\cis(10\degree)^4 \cdot \cis(16\degree)^5 $\\

$ = 2^4 \operatorname{name}{cis}(10\degree \cdot 4)  \cdot 1^5 \operatorname{cis}(16\degree \cdot 5)$\\

$ = 16\operatorname{cis}(40\degree) \cdot \operatorname{cis}(80\degree) $\\

$ = 16 \cis(120\degree) $\\ 

$ x = 16\cos \left(\dfrac{2\pi}{3}\right) = 16 \cdot \dfrac{-1}{2} = -8 $\\

$ y = 16\sin\left(\dfrac{2\pi}{3}\right) = 16 \cdot \dfrac{\sqrt{3}}{2} = 8\sqrt{3} $\\

$$ \text{Antwoord: } -8 + 8\sqrt{3}i $$\\

\begin{thebibliography}{999}
%\bibitem{a} Referentie
\end{thebibliography}
\end{document} 
